% 图 81.2 —— 由 `checks/emit_hand.py` 自动拼出（probe 精测 + prims2 图元分解）
%%
% ⚠ 这是**稿子不是成品**：跑 `checks/checkhand.py 81.2` 看指标 + 看叠合图。
%%
%% ⚠ 待核：
%   · 本图 probe 报出阴影族 1 个 —— **本脚本不生成阴影**（要照原书手写）；prims2 的 strokes 里可能含阴影线，核对时留意别画重
%   · 可疑字形 1 项（空心圈 ⊙ 常被认成字母 o）—— 见文件末
%%
\documentclass[12pt]{standalone}
\usepackage{fontspec}
\setsansfont{Arial}
\newfontfamily\mfallback{DejaVu Sans}
\usepackage{tikz}
\begin{document}
\begin{tikzpicture}[x=1pt, y=-1pt, line width=4.4pt, line cap=round, line join=round]

  %% ════ 填充区（prims2 的 fills）════

  %% ════ 直线段（prims2 的 strokes）════
  \draw[line width=7.7pt] (170.0,196.0) -- (172.0,189.0);
  \draw[line width=7.0pt] (815.0,188.0) -- (812.0,194.0);
  \draw[line width=6.5pt] (772.0,514.0) -- (763.0,504.0);
  \draw[line width=7.5pt] (721.0,348.0) -- (725.0,342.0);
  \draw[line width=7.0pt] (837.0,339.0) -- (798.0,339.0);
  \draw[line width=7.3pt] (761.0,504.0) -- (754.0,505.0);
  \draw[line width=6.7pt] (608.0,187.0) -- (606.0,193.0);
  \draw[line width=6.0pt] (389.0,435.0) -- (394.0,432.0);
  \draw[line width=3.0pt] (164.0,340.0) -- (161.8,342.0);
  \draw[line width=5.0pt] (161.8,342.0) -- (19.0,343.0);
  \draw[line width=5.0pt] (838.0,340.0) -- (946.0,341.0);
  \draw[line width=4.0pt] (396.0,247.0) -- (397.0,151.0);
  \draw[line width=7.0pt] (631.0,339.0) -- (591.0,339.0);
  \draw[line width=3.0pt] (631.0,339.0) -- (633.2,341.0);
  \draw[line width=5.0pt] (633.2,341.0) -- (724.0,341.0);
  \draw[line width=5.3pt] (1068.0,552.0) -- (1073.0,551.0);
  \draw[line width=5.0pt] (494.0,473.0) -- (488.0,477.0);
  \draw[line width=4.0pt] (396.0,526.0) -- (395.0,433.0);
  \draw[line width=5.0pt] (522.0,76.0) -- (528.0,73.0);
  \draw[line width=4.0pt] (396.0,341.0) -- (396.0,250.0);
  \draw[line width=6.1pt] (611.0,199.0) -- (607.0,195.0);
  \draw[line width=5.0pt] (397.0,342.0) -- (506.0,341.0);
  \draw[line width=3.0pt] (590.0,339.0) -- (587.8,341.0);
  \draw[line width=5.0pt] (587.8,341.0) -- (509.0,342.0);
  \draw[line width=4.0pt] (395.0,572.0) -- (394.0,595.0);
  \draw[line width=7.4pt] (389.0,252.0) -- (395.0,248.0);
  \draw[line width=7.0pt] (165.0,340.0) -- (203.0,340.0);
  \draw[line width=6.1pt] (402.0,252.0) -- (397.0,249.0);
  \draw[line width=3.0pt] (797.0,339.0) -- (794.8,341.0);
  \draw[line width=4.0pt] (794.8,341.0) -- (726.0,341.0);
  \draw[line width=5.0pt] (395.0,343.0) -- (395.0,431.0);
  \draw[line width=4.0pt] (681.0,420.0) -- (761.0,501.0);
  \draw[line width=6.0pt] (725.0,340.0) -- (722.0,335.0);
  \draw[line width=5.0pt] (394.0,342.0) -- (295.0,342.0);
  \draw[line width=4.0pt] (293.0,342.0) -- (205.2,342.0);
  \draw[line width=3.0pt] (205.2,342.0) -- (203.0,340.0);
  \draw[line width=4.0pt] (396.0,19.0) -- (398.0,108.0);
  \draw[line width=6.1pt] (396.0,432.0) -- (401.0,435.0);

  %% ════ 圆 / 椭圆（probe 精测）════
  \draw (611.4,266.4) circle (72.1);
  \draw (866.8,596.5) circle (72.3);
  \draw (184.5,267.0) circle (72.1);
  \draw (467.7,548.8) circle (72.5);
  \draw (470.2,129.0) circle (72.3);
  \draw (1009.8,597.2) circle (72.1);
  \draw (817.8,266.7) circle (72.1);

  %% ════ 实心点 ════
  \fill (517.5,80.5) circle (5.1);
  \fill (814.0,198.5) circle (5.0);
  \fill (294.5,338.5) circle (4.5);
  \fill (295.8,347.5) circle (5.4);
  \fill (508.5,347.5) circle (4.7);
  \fill (766.0,500.0) circle (5.8);
  \fill (791.0,588.3) circle (4.8);

  %% ════ 标签（probe 直接给的 \node 行）════
  %% ⚠ 全部补了 fill=white：文字在最上层，否则与曲线交叉干扰
     \node[anchor=center,inner sep=0pt,fill=white,font=\fontsize{30.3}{37.9}\selectfont] at (557.8,62.5) {\textsf{\textbf{\textit{A}}}};   % box(545,51,20,22)
     \node[anchor=center,inner sep=0pt,fill=white,font=\fontsize{22.2}{27.8}\selectfont] at (573.1,77.4) {\textsf{\textbf{1}}};   % box(567,69,7,16)
     \node[anchor=center,inner sep=0pt,fill=white,font=\fontsize{31.0}{38.8}\selectfont] at (220.8,178.0) {\textsf{\textbf{\textit{B}}}};   % box(209,166,20,23)
     \node[anchor=center,inner sep=0pt,fill=white,font=\fontsize{30.3}{37.9}\selectfont] at (655.8,178.5) {\textsf{\textbf{\textit{B}}}};   % box(644,167,20,22)
     \node[anchor=center,inner sep=0pt,fill=white,font=\fontsize{30.3}{37.9}\selectfont] at (865.8,178.5) {\textsf{\textbf{\textit{B}}}};   % box(854,167,20,22)
     \node[anchor=center,inner sep=0pt,fill=white,font=\fontsize{20.6}{25.7}\selectfont] at (249.4,192.3) {\textsf{\textbf{1}}};   % box(244,184,6,16)
     \node[anchor=center,inner sep=0pt,fill=white,font=\fontsize{20.6}{25.7}\selectfont] at (671.4,192.3) {\textsf{\textbf{1}}};   % box(666,184,6,16)
     \node[anchor=center,inner sep=0pt,fill=white,font=\fontsize{22.9}{28.7}\selectfont] at (880.1,192.4) {\textsf{\textbf{2}}};   % box(874,184,11,16)
     \node[anchor=center,inner sep=0pt,fill=white,font=\fontsize{22.1}{27.6}\selectfont] at (240.3,195.8) {{\mfallback\textbf{−}}};   % box(229,191,12,3)
     \node[anchor=center,inner sep=0pt,fill=white,font=\fontsize{31.4}{39.3}\selectfont] at (417.8,366.1) {\textsf{\textbf{\textit{e}}}};   % box(408,356,16,17)
     \node[anchor=center,inner sep=0pt,fill=white,font=\fontsize{24.9}{31.1}\selectfont] at (431.1,375.7) {\textsf{\textbf{0}}};   % box(423,367,12,17)
     \node[anchor=center,inner sep=0pt,fill=white,font=\fontsize{29.8}{37.3}\selectfont] at (761.0,437.5) {\textsf{\textbf{\textit{P}}}};   % box(750,425,18,23)
     \node[anchor=center,inner sep=0pt,fill=white,font=\fontsize{29.6}{37.0}\selectfont] at (532.2,463.5) {\textsf{\textbf{\textit{A}}}};   % box(520,452,19,22)
     \node[anchor=center,inner sep=0pt,fill=white,font=\fontsize{23.7}{29.7}\selectfont] at (561.8,477.4) {\textsf{\textbf{1}}};   % box(555,469,8,16)
     \node[anchor=center,inner sep=0pt,fill=white,font=\fontsize{22.1}{27.6}\selectfont] at (552.3,480.8) {{\mfallback\textbf{−}}};   % box(541,476,12,3)
     \node[anchor=center,inner sep=0pt,fill=white,font=\fontsize{29.6}{37.0}\selectfont] at (1112.2,565.5) {\textsf{\textbf{\textit{A}}}};   % box(1100,554,19,22)
     \node[anchor=center,inner sep=0pt,fill=white,font=\fontsize{30.7}{38.3}\selectfont] at (964.8,603.8) {\textsf{\textbf{\textit{b}}}};   % box(955,591,17,23)
     \node[anchor=center,inner sep=0pt,fill=white,font=\fontsize{31.0}{38.8}\selectfont] at (762.8,609.0) {\textsf{\textbf{\textit{B}}}};   % box(751,597,20,23)
     \node[anchor=center,inner sep=0pt,fill=white,font=\fontsize{24.9}{31.1}\selectfont] at (979.1,616.7) {\textsf{\textbf{0}}};   % box(971,608,12,17)

  % 钉死画布：否则超出的图元/控制点会把页面撑大，1:1 回比就失准
  \pgfresetboundingbox
  \path[use as bounding box] (0,0) rectangle (1129,682);
\end{tikzpicture}
\end{document}

% ⚠ 待核清单（probe 拿不准，人眼过一遍）：
%   · 未识别/跳过：⑤ 跳过/未识别的字形
